\section{假说：Long CoT 中的稳定“分子结构”}

为理解上述现象，我们借助图~\ref{fig:main} 与图~\ref{fig:intro} 的直观示例：高效 Long CoT 会在逻辑空间中把节点通过不同推理行为连接起来，从全局看形成具有互相支撑部件的稳定大分子结构。

我们把 Long CoT 轨迹形式化为一个带行为导向的图 \(G = (V, E)\)。其中，每个节点 \(v \in V\) 表示一个推理步骤，每条边 \(s = (u, v) \in E\) 被赋予一个行为标签 \(s \rightarrow b \in \{\mathcal{D}, \mathcal{R}, \mathcal{E}\}\)。对一个轨迹语料 \(\mathcal{C}\) 而言，我们估计其行为转移 \(P_{\mathcal{C}}(b' \mid b)\) 与边的边际分布 \(\pi_{\mathcal{C}}(b)\)。经验上，强推理教师在不同模型与任务上都会产生稳定的 \((P_{\mathcal{C}}, \pi_{\mathcal{C}})\)。尤其包括以下三类主要“键”\footnote{形式化定义与证明见附录~\ref{append:math}。}：\vspace{-8pt}







\paragraph{\textbf{深层推理如共价键}}
深层推理构成思维骨架，类似于共价键确定分子的主链。它编码强逻辑依赖（“步骤 A 必须支撑步骤 B”），维持方向性与连续性；一旦该骨架被破坏，后续步骤就会动摇，最终答案难以稳定。相比之下，“常规操作”更像每个步骤内部的局部稳定键，负责显式计算或直接语义表达。\vspace{-8pt}

\paragraph{\textbf{自我反思如氢键}}
反思是重要的稳定器。正如蛋白质链通过折叠并形成链内氢键来增强稳定，推理在后续步骤（如第 100 步）对早先前提（如第 10 步）进行检验、修正或强化时，也会获得稳定性。这种远程链接抑制了漂移与幻觉，使得长序列可以折叠成更自洽的结构。如果后续检查无法与先前承诺对齐，则意味着推理无法折叠，暴露结构性逻辑错误。

\paragraph{\textbf{自我探索如范德华力}}
探索类似短暂的范德华作用：它支撑溯因或归纳式的扩展，通过在语义空间里建立低承诺的关联，让概念得以暂时游移、组合、验证，再由更强约束收敛。
